% Trigonometrie – bank/trigonometrie/ (zusammenbau v0.4, Heft)
\einheitenkopf[t5]{Trigonometrie}
\setcounter{aufgabe}{4}

\begin{aufgabe}{Seite berechnen – Prüfungsaufgaben}
\begin{teile}
\teil Rechter Winkel bei $C$. Trage den Bruch ein. \hfill (P10 2017 OS) \\

\dreieck{(0,0)}{(5,0)}{(1.8,2.4)}{d}{e}{f}{}{\beta}{}
$\mathrm{sin}\,\beta =$ \leerfeld
\teil Rechter Winkel bei $B$. Trage den Bruch ein. \hfill (P10 2019 OS) \\

\dreieck{(0,0)}{(5,0)}{(5,3)}{k}{m}{n}{\alpha}{}{}
$\mathrm{cos}\,\alpha =$ \leerfeld
\teil Rechter Winkel bei $A$. Welche Gleichung gilt? \hfill (P10 2018 OS) \\ \kreuz{$\mathrm{sin}\,\beta = \frac{p}{q}$} \\ \kreuz{$\mathrm{cos}\,\beta = \frac{q}{p}$} \\ \kreuz{$\mathrm{sin}\,\beta = \frac{q}{p}$}

\dreieck{(0,0)}{(4,0)}{(0,3)}{p}{q}{r}{}{\beta}{}
\teil Rechter Winkel bei $C$. Welche Gleichung gilt? \hfill (P10 2018 OS) \\ \kreuz{$\mathrm{cos}\,\alpha = \frac{w}{v}$} \\ \kreuz{$\mathrm{cos}\,\alpha = \frac{v}{w}$} \\ \kreuz{$\mathrm{sin}\,\alpha = \frac{v}{w}$}

\dreieck{(0,0)}{(5,0)}{(3.2,2.4)}{u}{v}{w}{\alpha}{}{}
\teil Rechter Winkel bei $B$. Welche Gleichung gilt? \hfill (P10 2018 OS) \\ \kreuz{$\mathrm{tan}\,\gamma = \frac{m}{k}$} \\ \kreuz{$\mathrm{tan}\,\gamma = \frac{k}{m}$} \\ \kreuz{$\mathrm{sin}\,\gamma = \frac{m}{n}$}

\dreieck{(0,0)}{(5,0)}{(5,3)}{m}{n}{k}{}{}{\gamma}
\end{teile}
\end{aufgabe}

\begin{aufgabe}{Seite berechnen – Prüfungsaufgaben – weiter}
\begin{teile}
\teil $\mathrm{cos}\,60^\circ = \frac{4}{x}$ – berechne $x$. \hfill (P10 2020 OS) \\ x = \leerfeld
\teil $\mathrm{tan}\,45^\circ = \frac{6}{x}$ – berechne $x$. \hfill (P10 2020 OS) \\ x = \leerfeld
\teil Rechter Winkel bei $B$. Trage den Bruch ein. \hfill (P10 2025 OS) \\

\dreieck{(0,0)}{(5,0)}{(5,3)}{e}{f}{d}{}{}{\gamma}
$\mathrm{sin}\,\gamma =$ \leerfeld
\teil Rechter Winkel bei $A$. Trage den Bruch ein. \hfill (P10 2025 OS) \\

\dreieck{(1.8,2.4)}{(0,0)}{(5,0)}{z}{x}{y}{}{\beta}{}
$\mathrm{sin}\,\beta =$ \leerfeld
\teil Rechter Winkel bei $C$. Gib $\mathrm{tan}\,\alpha$ als Bruch an. \hfill (P10 2020 OS) \\

\dreieck{(4,0)}{(0,3)}{(0,0)}{q}{p}{s}{\alpha}{}{}
$\mathrm{tan}\,\alpha =$ \leerfeld
\teil Rechter Winkel bei $A$. Gib $\mathrm{tan}\,\gamma$ als Bruch an. \hfill (P10 2020 OS) \\

\dreieck{(0,0)}{(0,3)}{(5,0)}{w}{v}{u}{}{}{\gamma}
$\mathrm{tan}\,\gamma =$ \leerfeld
\end{teile}
\end{aufgabe}
